
\thispagestyle{plain}

% -----------------------------------------
% TITRE
% -----------------------------------------

\vspace*{12mm}

\begin{center}
%    {\titlefont\fontsize{30}{36}\selectfont
    Alban : montagne druidique !
%}
\end{center}

\vspace{23mm}

% -----------------------------------------
% PREMIER PARAGRAPHE
% -----------------------------------------

{\itshape

Q{uelle}
que soit l'origine du terme « Alban », il est certain
qu'à l'âge de la pierre polie (au néolithique) quelque
2000 ans avant Jésus Christ, des hommes vécurent
sur notre petit plateau et y dressèrent le menhir et le
dolmen qu'on y voit encore. Ceux-ci portent respectivement
le nom de pierre du diable et palet de Notre Dame.
La légende mariale attachée à ces mégalithes est très
pittoresque et fera l'objet d'une relation.

}

\vspace{13mm}

% -----------------------------------------
% DEUXIEME PARAGRAPHE
% -----------------------------------------

{\itshape

D{es}
découvertes faites dans les chantiers d'exploitation
de la région minière qui s'étend d'Ambialet au Roc
de Saint Michel ont établi que des gisements ferreux
étaient exploités dès l'antiquité la plus haute.
Des haches, des lances, des bracelets trouvés nous
révèlent que 12 siècles avant le Christ, les Ligures
parcouraient notre région. Plus tard, les Celtes,
800 ans avant J.C. exploiteront la région albanaise.

}
